\documentclass[10pt]{article}

\usepackage[utf8]{inputenc}

\usepackage[T1]{fontenc}

\usepackage{amsmath}

\usepackage{amsfonts}

\usepackage{amssymb}

\usepackage[version=4]{mhchem}

\usepackage{stmaryrd}



\begin{document}

Уравнение Фоккера - Планка для переходной плотности вероятности О.- У. п. имеет вид:



\[

\frac{\partial}{\partial t} p\left(u t \mid u_{0} t_{0}\right)=\frac{\partial}{\partial u} \gamma u p\left(u t \mid u_{0} t_{0}\right)+\frac{1}{2} \sigma \frac{\partial^{2}}{\partial u^{2}} p\left(u t \mid u_{0} t_{0}\right),

\]



фундаментальное решение к-рого есть



\[

\begin{aligned}

& p\left(u t \mid u_{0} t_{0}\right)=\left(\frac{\pi \sigma}{\gamma}\left(1-e^{-2 \gamma\left(t-t_{0}\right)}\right)\right)^{-1 / 2} \times \\

& \times \exp \left\{-\frac{\left(u-u_{0} e^{-\gamma\left(t-t_{0}\right)}\right)^{2}}{\frac{\sigma}{\gamma}\left(1-e^{-2 \gamma\left(t-t_{0}\right)}\right)}\right\}, \quad t>t_{0} .

\end{aligned}

\]



О.- У. п. является гауссовским стационарным процессом и обладает следуюшими свойствами:



a) $\mathrm{E}\left\{U_{t}\right\}=0, \mathrm{D}\left\{U_{t}\right\}=\frac{\sigma}{2 \gamma}\left(1-e^{-2 \gamma t}\right), \mathrm{E}\left\{U_{t} U_{s}\right\}=\frac{\sigma}{2 \gamma} e^{-\gamma|t-s|}$, причем при вычислении этих средних следует выполнить усреднение по начальным данным $u_{0}$ O.- У. п. со стационарной функцией распределения



\[

p\left(u_{0}\right)=\left(\frac{\pi \sigma}{\gamma}\right)^{-1 / 2} \exp \left\{-\frac{u_{0}^{2}}{\sigma / \gamma}\right\}

\]



б) О.- У. П. не имеет независимых приращений, то есть изменения скорости броуновской частицы коррелированы на непересекающихся временных интервалах.



О.- У. п. является единственным стационарным, гауссовским, марковским случайным процессом (те орема Дуба).



Стационарность О.- У. п. и отсутствие независимых приращений отличают его от винеровского процесса. Винеровский процесс, как математич. модель броуновского движения в пространстве координат блуждающей частицы, получается из О.-У. п. в результате предельного перехода $\sigma \rightarrow \infty$, $\gamma \rightarrow \infty$, так что $\frac{\sigma}{\gamma^{2}}=$ const (для стандартного винеровского процесса $\frac{\sigma}{\gamma^{2}}=1$ ). Таким образом, винеровский процесс соответствует описанию броуновского движения в пределе большой вязкости среды (сильное трение) и интенсивного шума.



О.- У.п. $U_{\tau}$ при $0 \leqslant \tau \leqslant t$ допускает канон ическое представление:



\[

U_{\tau}=\sum_{k=0}^{\infty} u_{k} v_{k}(\tau)

\]



где коэффициенты $u_{k}$ - независимые гауссовские случайные величины такие, что $\mathrm{E}\left(u_{k}\right)=0, \mathrm{E}\left(u_{k} u_{l}\right)=\lambda_{k} \delta_{k l}, \lambda_{k}-$ собственные числа, $v_{k}(\tau)$ - собственные функции интегрального оператора с ядром $\mathrm{E}\left\{U_{\tau} U_{s}\right\}=\frac{\sigma}{2 \gamma} e^{-\gamma|\tau-s|}$.



\[

\int_{0}^{t} d s \mathrm{E}\left\{U_{\tau} U_{s}\right\} v_{k}(s)=\lambda_{k} v_{k}(\tau)

\]



М ногоме р ный О.- У. п. $U_{t}=\left\{U_{t}^{(1)}, \ldots, U_{t}^{(n)}\right\}$ может быть определен как решение уравнения Ланжевена



\[

d U_{t}=-\Gamma U_{t} d t+\mathrm{B} d W_{t}

\]



где $W_{t}=\left\{W_{t}^{(1)}, \ldots, W_{t}^{(n)}\right\}, W_{t}^{(i)}-$ стохастически независимые стандартные винеровские процессы $\mathrm{E}\left\{W_{t}^{(i)} W_{s}^{(k)}\right\}=\delta_{i k} \min (t, s)$, $\|\Gamma\|_{i k}$ и $\|\mathrm{B}\|_{i k}$ суть $(n \times n)$-матрицы, причем матрица $\|\mathrm{B}\|_{i k}$ связана с диффузионной матрицей $\|\sigma\|_{i k}$, входящей в уравнение Фоккера - Планка соотношением $\|\sigma\|_{i k}=\left\|\mathrm{BB}^{\top}\right\|_{i k}$. Уравнение Фоккера - Планка для многомерного О.- У. п. имеет вид:



$\frac{\partial}{\partial t} p\left(\boldsymbol{u} t \mid \boldsymbol{u}_{0} t_{0}\right)=\frac{\partial}{\partial u_{i}} \Gamma_{i k} u_{k} p\left(\boldsymbol{u} t \mid \boldsymbol{u}_{0} t_{0}\right)+\frac{1}{2} \sigma_{i k} \frac{\partial^{2}}{\partial u_{i} \partial u_{k}} p\left(\boldsymbol{u} t \mid \boldsymbol{u}_{0} t_{0}\right)$,



\section{OРНШТЕЙНА}

фундаментальное решение к-рого есть



\[

p\left(\boldsymbol{u} t \mid \boldsymbol{u}_{0} t_{0}\right)=(2 \pi)^{-n / 2}(\operatorname{det} A)^{-1 / 2} \exp \left\{-\tilde{u}_{i}\left(A^{-1}\right)_{i k} \tilde{u}_{k}\right\},

\]



где



\[

\begin{gathered}

\tilde{u}_{i}=u_{i}-\left(e^{-\Gamma\left(t-t_{0}\right)}\right)_{i k} u_{0 k}, \\

A_{i k}=\int_{t_{0}}^{t} d \tau\left(e^{-\Gamma(t-\tau)}\right)_{i j} \sigma_{j l}\left(e^{-\Gamma(t-\tau)}\right)_{l k} .

\end{gathered}

\]



Простейшим примером двумерного $(n=2)$ O.- У.п. может служить гармонический стохастич. процесс $x_{t}$, описываемый стохастич. дифференциальным уравнением



\[

d \dot{x}_{t}+2 \beta d x_{t}+\omega_{0} x_{t} d t=\sqrt{\sigma} d W_{t},

\]



где $W_{t}$ - стандартный винеровский процесс, $\omega_{0}$ - частота свободных колебаний, $\beta-$ коэффициент затухания, $\sigma-$ интенсивность внешнего шума. В терминах двухкомпонентных векторов $U_{t}=\left\|\begin{array}{l}x_{t} \\ \dot{x}_{t} \\ \omega\end{array}\right\|$ и $d w_{t}=\left\|\begin{array}{l}0 \\ \frac{\sqrt{\sigma}}{\omega_{0}} d W_{t}\end{array}\right\|$ уравнение Ланжевена для стохастич. гармонич. процесса приобретает вид



\[

d \boldsymbol{U}_{t}=-\Gamma \boldsymbol{U}_{t} d t+d W_{t}

\]



где матрица $\Gamma=\left\|\begin{array}{cr}0 & -\omega_{0} \\ \omega_{0} & 2 \beta\end{array}\right\|$, а корреляционная функция двумерного процесса $\boldsymbol{W}_{\boldsymbol{t}}$ имеет вид:



\[

\mathrm{E}\left\{W_{t}^{(i)} W_{s}^{(k)}\right\}=\tilde{\sigma}_{i k} \min (t, s),

\]



где матрица $\tilde{\sigma}=\frac{\sigma}{\omega_{0}}\left\|\begin{array}{ll}0 & 0 \\ 0 & 1\end{array}\right\|$.



Jum.: [1] Uhlenbeck G. E., Ornste in L. S., «Phys. Rev.», 1930, v. 36, p. 823-41; [2] Ч андр а с е ка р С., Стохастические проблемы в физике и астрономии, пер. с англ., М., 1947; [3] Ги х ман И. И., Ско ро х од А. В., Теория случайных процессов, т. 2., М., 1972.



Н. В. Ласкин.



OPHШТЕЙНА - ЦЕРНИКЕ УРАВНЕНИЕ - ТОчНОе СОотношение, связывающее парную корреляиионную функцию $\mathrm{v}_{2}(r)$ и прямую корреляционную функцию $C(r)$ пространственно однородной системы:



\[

v_{2}\left(\overrightarrow{r_{1}}-\overrightarrow{r_{2}}\right)=C\left(\overrightarrow{r_{1}}-\overrightarrow{r_{2}}\right)+n \int d \overrightarrow{r_{3}} v_{2}\left(\overrightarrow{r_{1}}-\overrightarrow{r_{3}}\right) C\left(\overrightarrow{r_{3}}-\overrightarrow{r_{2}}\right)

\]



где $n-$ средняя плотность числа частиц в системе. О.- Ц. у. содержит две неизвестные функции. Чтобы дополнить О.Ц. у., используются приближенные уравнения различной степени сложности (см. Перкуса - Иевика уравнение, Гиперцепное уравнение).



Jum.: [1] Ornstein L. S., Zernike F., «Proc. Acad. Sci. Amst.», 1914, t. 17, р. 783; [2] Б але с ку Р., Равновесная и неравновесная статистическая механика, пер. с англ., т. 1, М., 1978.



Д. П. Санкович. ОРРА-ЗОММЕРФЕЛЬДА УРАВНЕНИЕ - ДИФФеРеНЦИальное уравнение



\[

\varphi^{\mathrm{IV}}-2 \alpha^{2} \varphi^{\prime \prime}+\alpha^{4} \varphi=i \alpha \operatorname{Re}\left\{(\omega-c)\left(\varphi^{\prime \prime}-\alpha^{2} \varphi\right)-\omega^{\prime \prime} \varphi\right\}

\]



для амплитуды двумерных возмушений плоскопараллельного течения вязкой жидкости, имеющих вид распространяющихся волн



\[

\varphi(y) \exp (i \alpha(x-c t)) \text {. }

\]



Здесь $\operatorname{Re}$ - число Рейнольдса течения, $w(y)$ - профиль скорости, $c$ - собственное значение, подлежащее определению. Обычно на плоскости параметров $\alpha$, Re строится кривая нейтральной устойчивости, соответствующая $\operatorname{Im} c=0$. Для плоского течения Пуазейля установлено наличие неустойчивости. Теория численного интегрирования О.-3. у. построена с использованием ВКБ-асимптотик (см.[3]).



Jum.: [1] O rr W., «Proc. Roy Irish Acad.», A, 1906-07, v. 27, p. 927, 69-138; [2] So m me rfeld A., «Atti del IV Congr. intem. dei matem. Roma», 1909, v. 3, p. 116-24; [3] Л и н ь Цзя - цзя о, Теория гидродинамической устойчивости, пер. с англ., М., 1958. Л. А. Дикий.





\end{document}